\section{Related work}\label{sec:related}

\paragraph{Petersen colorings and normal colorings.} The Petersen coloring conjecture is Jaeger's
\cite{Jaeger1988}; the equivalence with normal 5-edge-colorings is his~\cite{Jaeger1985}. Normal
edge-colorings of cubic graphs were studied by Mazzuoccolo and Mkrtchyan~\cite{MM2020}, who also stated,
for the normal chromatic index, the conjecture that it is at most 6~\cite{MM2020b}. The paper that
disproved the Petersen coloring conjecture on 52 vertices~\cite{GJMMMU2026} presents Petersen colorings
as cycle-continuous maps into the Petersen graph and as zero-temperature states of spin models, and it
follows the first announcement of the counterexamples~\cite{GJMMM2026zen}. The charges used below make
the second reading concrete: a vertex whose three labels do not form a star carries a nonzero class
modulo the cut space, and a Petersen coloring is a configuration without charge. The Petersen coloring
conjecture states that the Petersen graph alone colors every bridgeless cubic graph; the sets of
$r$-graphs that color all $r$-graphs are studied in~\cite{MMSW2025}, and a 2026 search for coloring
results on $r$-graphs found~\cite{MSWZ2024} on perfect-matching conjectures, without a statement on
Petersen colorings. Colorings of cubic graphs by the Sylvester graphs, a neighbouring family of
$H$-colorings, are the subject of~\cite{Ferrarini2026}\,\claim{U}. The paper~\cite{GJMMMU2026} also cites~\cite{Inoue2026}, on 3-edge-colorings of apex cubic graphs, as
completing the framework of Tutte's conjecture that bridgeless cubic graphs without a Petersen minor are 3-edge-colorable\,\claim{U}.

\paragraph{Abnormal edges and sublinear bounds.} The least number of abnormal edges is the quantity of
Mattiolo, Mazzuoccolo and Mkrtchyan~\cite{MMM2021}, who reduced the sublinear bound for bridgeless
graphs to the Petersen coloring conjecture and posed Conjecture~\ref{conj:mmm}. After the
disproof, rings and frames of counterexamples settled the classes joined through 2-edge and 3-edge
cuts~\cite{Audit2026}; the present paper settles the last class. A 2026 search of arXiv for work on
abnormal edges returned only~\cite{MMM2021}.

\paragraph{Multipoles, snark families and superpositions.} The flower snarks of Isaacs~\cite{Isaacs1975}
are rings of claws, and H\"agglund and Steffen showed that they have Petersen colorings and studied
further families of snarks~\cite{HS2014}. Sedlar and \v{S}krekovski study superpositions of snarks by
the Petersen graph, whose basic piece is the superedge used in
Section~\ref{sec:theory6}~\cite{SS2023}. Generation and properties of snarks are the subject
of~\cite{Brinkmann2013}, from whose censuses the cyclically 5-edge-connected candidates for the question
of Section~\ref{sec:questions} would be taken; the smallest snarks with oddness 4 and cyclic
connectivity 4 are studied in~\cite{GMS2019}\,\claim{U}. The counterexamples on 52 vertices and in every even order from 60 are those
of~\cite{GJMMMU2026}; the edge list of the first one is that of its House of Graphs entry~\cite{HoG2023}.

\paragraph{Cut-space and transfer methods.} The argument in Section~\ref{sec:classes}
is a cut-space argument: it uses only the cut code of the Petersen graph, a binary $[15,9,3]$ code, and
the fact that a vertex sum of stars is a cut. Both papers of Jaeger~\cite{Jaeger1988,Jaeger1985} belong to
the theory of nowhere-zero flows, a setting of the same kind. We did not find, in the sources we hold, a treatment of
the line-graph distance as the complete invariant of a pair of labels modulo the cut space, nor an earlier use of
transfer relations over cut-space sectors with a finite-semigroup certificate for rings of multipoles of
snark families (\claim{U}: the search covered the sources listed in the references and the 2026 arXiv listings
for the Petersen coloring conjecture and abnormal edges, not the whole literature on transfer matrices).
The questions that follow from this gap are stated in Section~\ref{sec:questions}.

\paragraph{Companion results.} The two 52-vertex counterexamples are colorable only by themselves, if no
counterexample has at most 38 vertices~\cite{COBI2026}; the constructions below use the first of them and
do not depend on that result. The consequences of the counterexamples for covers, flows and normality
defects are in~\cite{Audit2026}. The open problems of the conjecture are listed in~\cite{OPG}.
